\subsection{Identification des goulots d'étranglement --- itération~2}
\label{sec:scenario1-profil2}

Après application des trois premières optimisations, une seconde session
de profilage a été conduite selon le même protocole que la
section~\ref{sec:scenario1-profil1}. Les figures~\ref{fig:jfr-2}
et~\ref{fig:prof2-zoom} avaient déjà confirmé l'efficacité des
corrections précédentes~: \texttt{findMetaObjectsInDefinition\_oneshot}
n'apparaît plus de manière dominante dans le profil CPU. Un nouveau
goulot d'étranglement émerge cependant, centré cette fois sur
\texttt{\_isInDefinition} et ses appels à \texttt{getDomain}.

% -----------------------------------------------
\subsubsection{Profilage CPU}
% -----------------------------------------------

En se concentrant sur les méthodes les plus représentées dans le
flamegraph post-optimisation~1--3, on obtient les pourcentages d'occupation
CPU suivants lors du bootstrapping~:

\begin{itemize}
    \item \texttt{pcGetdeleted} --- 21,33\%
    \item \texttt{check} --- 26,51\%
    \item \texttt{getDomain} --- 37,59\%
    \item \texttt{\_isInDefinition} / \texttt{isInDefinition} --- 37,23\%
    \item \texttt{\_getName} --- 48,43\%
    \item \texttt{getMOByName} --- 50,48\%
\end{itemize}

\texttt{getDomain} et \texttt{\_isInDefinition} apparaissent au même
niveau de profondeur dans le flamegraph, ce qui indique que
\texttt{getDomain} est appelé massivement \textit{depuis}
\texttt{\_isInDefinition}. La nouvelle chaîne d'appels critique est~:
\begin{center}
\small
\texttt{getMOByName} \\
$\downarrow$ \\
\texttt{\_getName}
(\texttt{MetaObject\_impl},
\texttt{metadone\_bundle\_metabusiness}) \\
$\downarrow$ \\
\texttt{isInDefinition} $\to$ \texttt{\_isInDefinition} 
(\texttt{ConcreteModel\_impl},
\texttt{metadone\_bundle\_metabusiness}) \\
$\downarrow$ \\
\texttt{getDomain} 
(interface \texttt{K\_Property},
\texttt{metadone\_bundle\_repository})
\end{center}

Notons que \texttt{getMOByName} et \texttt{\_getName} occupent toujours
le haut du profil (48--50\%) non plus à cause de la précondition
éliminée par l'optimisation~2, mais à cause du volume d'appels légitime
qu'ils génèrent vers \texttt{isInDefinition}. Le goulot a donc migré
d'un niveau vers le bas dans la pile d'appels.

% -----------------------------------------------
\subsubsection{Analyse statique et quantification}
% -----------------------------------------------

%\paragraph{Rôle de \texttt{\_isInDefinition}.}
La méthode \texttt{\_isInDefinition} est définie dans
\texttt{ConcreteModel\_impl} du bundle
\texttt{metadone\_bundle\_metabusiness}. Elle répond à une question
symétrique de \texttt{findMetaObjectsInDefinition\_oneshot}~: un
\texttt{ConcreteObject} donné appartient-il à la définition de ce
\texttt{ConcreteModel}~?

Pour le comprendre, il faut rappeler comment l'appartenance d'un objet
concret à un modèle est encodée dans le kernel: de même qu'un \texttt{MetaDef} relie un
\texttt{MetaObject} à son \texttt{MetaModel} via une \texttt{K\_Property}
de type \texttt{def\_pt}, un \texttt{ConcreteDef} relie un
\texttt{ConcreteObject} à son \texttt{ConcreteModel} par la même
structure. La méthode itère donc sur l'ensemble des \texttt{K\_Property}
dont \texttt{this} est le range via \texttt{getTarget().getRangeOf()}, et
pour chaque candidat, vérifie deux conditions~: que son domaine est bien
l'objet recherché (\texttt{candidate.getDomain() == kmember}), et que son
type de propriété est \texttt{def\_pt} (\texttt{someParentIsDef}).
Il s'agit d'une \textbf{recherche linéaire $O(n)$}, où $n$ est le nombre
de propriétés dont le \texttt{ConcreteModel} est le range.

\begin{lstlisting}[
  caption={\texttt{\_isInDefinition} --- classe \texttt{ConcreteModel\_impl},
  bundle \texttt{metadone\_bundle\_metabusiness}. Recherche linéaire $O(n)$
  sur les \texttt{K\_Property} du range~: chaque itération appelle
  \texttt{getDomain()} (interface \texttt{K\_Property},
  bundle \texttt{metadone\_bundle\_repository}).},
  label={ls-isinDefinition}]
private boolean _isInDefinition(ConcreteObject_impl member) {
    try {
        final K_Object kmember = member.getTarget();
        for (final K_Property candidate :
                getTarget().getRangeOf()) {
            if (candidate.getDomain() == kmember
                    && someParentIsDef(candidate)) {
                return true;
            }
        }
        return false;
    } catch (final BadFacet e) {
        throw new Metal2Exception(e);
    } catch (final BadPreCondition e) {
        throw new Metal2Exception(e);
    }
}
\end{lstlisting}

%\paragraph{Source de l'amplification.}
Rappelons que \texttt{\_getName} (section~\ref{sec:opti-metamodel},
optimisation~2) parcourt les \texttt{ConcreteProperty<String>} associées
au \texttt{MetaObject} et appelle \texttt{concreteModel.isInDefinition(cp)}
pour chacune afin de déterminer si la propriété de nom appartient au bon
\texttt{ConcreteModel} (Listing~\ref{ls--getname}, ligne~378). Chaque
appel à \texttt{isInDefinition} déclenche un scan complet via
\texttt{\_isInDefinition}. Si la propriété n'est pas trouvée dans le
\texttt{MetaModel} courant, \texttt{\_getName} remonte récursivement dans
ses supertypes directs (ligne~383), multipliant encore le nombre d'appels.

L'instrumentation confirme cette amplification~: pour
\textbf{1~700~appels} à \texttt{\_getName}, on dénombre plus de
\textbf{74~000~appels} à \texttt{\_isInDefinition}, soit un
\textbf{facteur multiplicatif moyen de 43}. Ce ratio est du même ordre
que le facteur~57 observé à la section~\ref{sec:scenario1-profil1}~: la
structure du problème est identique --- une boucle externe génère une
amplification quadratique via une recherche linéaire interne --- mais le
point d'application s'est déplacé d'un niveau vers le bas dans la pile
d'appels, de \texttt{MetaModel\_impl} vers \texttt{ConcreteModel\_impl}.

%\paragraph{Parallèle avec le problème précédent.}
Le schéma identifié ici est la réplique exacte de celui de
l'itération~1 (section~\ref{sec:scenario1-profil1}), mais opérant
cette fois sur le niveau concret (M1) de MetaL$_2$ plutôt que sur
son niveau méta-modèle (M2/M3)~:

\begin{center}
\small
\begin{tabular}{lll}
\hline
\textbf{Dimension} & \textbf{Itération~1} & \textbf{Itération~2} \\
\hline
Niveau (dans MetaL$_2$) & Méta-modèle (M2/M3) & Concret (M1) \\
Classe & \texttt{MetaModel\_impl} & \texttt{ConcreteModel\_impl} \\
Méthode centrale & \texttt{\_hasInItsDefinition} &
\texttt{\_isInDefinition} \\
Scan interne & \texttt{findMetaObjectsInDefinition\_oneshot} &
\texttt{\_isInDefinition} elle-même \\
Opération de base & \texttt{getRangeOf(def\_pt)} & \texttt{getRangeOf()} \\
Facteur d'amplification & $\sim$57 & $\sim$43 \\
\hline
\end{tabular}
\end{center}

Cette récurrence est la conséquence directe de la symétrie réflexive
entre les niveaux méta et concret de MetaL$_2$, documentée à la
section~\ref{sec:reflexivite-intro}~: \texttt{MetaModel} et
\texttt{ConcreteModel} reposent tous deux sur le substrat uniforme
de \texttt{DataObject} (MetaL$_1$, bundle \texttt{repository}) et
encodent l'appartenance de la même manière. Les algorithmes qui la
vérifient partagent donc la même structure et les mêmes
vulnérabilités. La solution sera également symétrique~: l'introduction
d'un cache dans \texttt{ConcreteModel\_impl}, à l'image de ce qui a
été fait dans \texttt{MetaModel\_impl}.

% -----------------------------------------------
\subsubsection{Synthèse du diagnostic}
% -----------------------------------------------

Le goulot résiduel identifié après la première itération d'optimisation est
analogue à celui de l'itération~1, mais opère au niveau concret (M1) de
MetaL$_2$.
La pile d'appels problématique est~:

\begin{center}
\small
\texttt{ConcreteModel\_impl.\_isInDefinition} \\
(itération $O(n)$ sur toutes les \texttt{K\_Property} du range) \\[4pt]
$\downarrow$ \\[4pt]
\texttt{K\_Property.getDomain} ---
\texttt{metadone\_bundle\_repository} \\
(accès au dépôt JPA à chaque itération)
\end{center}

La mitigation naturelle est identique à celle de l'itération~1~: la
mémoïsation des résultats via un cache. Dans
\texttt{ConcreteModel\_impl}, l'appartenance d'un
\texttt{ConcreteObject} à un \texttt{ConcreteModel} est une information
qui change rarement une fois le modèle construit~; elle est un candidat
idéal pour la mise en cache. La phase d'optimisation suivante développe
cette stratégie.

\subsection{Optimisation --- itération~2}
\label{sec:scenario1-opti2}

\subsubsection{Optimisation~4 : cache d'appartenance dans
\texttt{ConcreteModel\_impl}}
\label{sec:opti4}

%\paragraph{Principe.}
Le diagnostic de la section~\ref{sec:scenario1-profil2} a mis en évidence
que \texttt{\_isInDefinition} était appelée 74~000 fois pour seulement
1~700~appels à \texttt{\_getName}, et que chaque appel déclenchait un scan
linéaire complet via \texttt{getRangeOf()} et \texttt{getDomain()} dans
le bundle \texttt{metadone\_bundle\_repository}. Le correctif introduit un cache
d'appartenance directement dans \texttt{ConcreteModel\_impl}
(\texttt{metadone\_bundle\_metabusiness}), afin que chaque paire
(\texttt{ConcreteModel}, \texttt{ConcreteObject}) ne soit scannée qu'une
seule fois.

Cette approche est la réplique exacte, au niveau concret (M1) de
MetaL$_2$, du mécanisme de \texttt{CacheCallFunction} utilisé dans
\texttt{MetaModel\_impl} pour \texttt{\_hasInItsDefinition}
(optimisations~1 et~3, section~\ref{sec:opti-metamodel}). La symétrie
réflexive entre les niveaux méta et concret de MetaL$_2$, documentée
à la section~\ref{sec:reflexivite-intro}, se retrouve ici dans la
solution autant que dans le problème.

%\paragraph{Structure du cache.}
Le cache est déclaré comme champ d'instance \texttt{private final} de
\texttt{ConcreteModel\_impl}~:

\begin{lstlisting}[firstnumber=42,
  caption={Déclaration du cache dans \texttt{ConcreteModel\_impl} ---
  bundle \texttt{metadone\_bundle\_metabusiness}.},
  label={ls-cache-isin}]
private final Map<ConcreteObject_impl, Boolean>
    cache_isInDefinition =
        new WeakHashMap<ConcreteObject_impl, Boolean>();
\end{lstlisting}

Le choix d'un \texttt{WeakHashMap} est délibéré et cohérent avec celui
fait pour le cache de proxies dans \texttt{FactoryOfProxy}~: les clés sont des références faibles vers
les \texttt{ConcreteObject\_impl}. Lorsqu'un objet n'est plus référencé
ailleurs dans la JVM, il peut être collecté par le ramasse-miettes sans
que le cache ne l'en empêche, et l'entrée correspondante disparaît
automatiquement. Cela évite toute fuite mémoire dans un système où les
objets concrets peuvent être supprimés à l'exécution via
\texttt{removeFromDefinition}.

Par ailleurs, le cache mémorise les \textbf{deux résultats possibles}
(\texttt{true} et \texttt{false})~: il évite ainsi de rescanner pour
des objets dont l'absence est déjà connue, ce qui est important dans les
appels de vérification défensifs effectués avant un ajout.

%\paragraph{Logique de consultation.}
La méthode \texttt{\_isInDefinition} réécrite est présentée au
Listing~\ref{ls-isinDefinition-cache}. Elle consulte désormais le
cache avant tout accès au repository~: en cas de succès
(\textit{cache hit}), la réponse est retournée en $O(1)$~; en cas
d'échec (\textit{cache miss}), le scan $O(n)$ est effectué et son
résultat --- positif ou négatif --- est mémorisé avant d'être
retourné.

\begin{lstlisting}[language=Java, firstnumber=141,
  caption={\texttt{\_isInDefinition} après l'optimisation~4 ---
  classe \texttt{ConcreteModel\_impl},
  bundle \texttt{metadone\_bundle\_metabusiness}. Le cache est
  consulté avant tout scan~; les deux résultats (positif et
  négatif) sont mémorisés.},
  label={ls-isinDefinition-cache}]
private boolean _isInDefinition(ConcreteObject_impl member) {
    try {
        // Check cache first
        Boolean cached = cache_isInDefinition.get(member);
        if (cached != null) {
            return cached;
        }

        // Cache miss: perform the scan
        final K_Object kmember = member.getTarget();
        for (final K_Property candidate : getTarget().getRangeOf()) {
            if (candidate.getDomain() == kmember
                    && someParentIsDef(candidate)) {
                cache_isInDefinition.put(member, true);
                return true;
            }
        }

        // Not found: store false and return
        cache_isInDefinition.put(member, false);
        return false;
    } catch (final BadFacet e) {
        throw new Metal2Exception(e);
    } catch (final BadPreCondition e) {
        throw new Metal2Exception(e);
    }
}
\end{lstlisting}

Notons que le gain de performance attendu ne porte pas sur le nombre
d'invocations de \texttt{\_isInDefinition} --- qui reste identique
puisque les call sites en amont (\texttt{\_getName} en particulier)
ne sont pas modifiés --- mais sur le \textbf{coût unitaire} de chaque
appel~: ce qui était auparavant un scan $O(n)$ du graphe devient,
dans la majorité des cas, un accès $O(1)$ à la \texttt{WeakHashMap}.
La mesure pertinente du gain est donc le temps mural, analysé à la
section~\ref{sec:scenario1-tests2}.

%\paragraph{Cohérence du cache lors des mutations.}
Deux méthodes modifient la définition d'un \texttt{ConcreteModel} et
doivent maintenir la cohérence du cache~:

\begin{itemize}
    \item \texttt{\_addInDefinition} insère directement \texttt{true}
    dans le cache après l'ajout. Cela évite qu'un appel immédiat à
    \texttt{\_isInDefinition} ne déclenche un scan inutile pour un objet
    dont on vient d'établir l'appartenance --- comportement analogue à la
    mise à jour incrémentale de l'optimisation~3
    (section~\ref{sec:opti-metamodel}).

    \item \texttt{removeFromDefinition} invalide l'entrée via
    \texttt{cache\_isInDefinition.remove(member)} après suppression. Au
    prochain appel, \texttt{\_isInDefinition} effectuera un scan et
    constatera l'absence, qu'elle mettra alors en cache avec
    \texttt{false}.
\end{itemize}

Le cache étant de portée instance, chaque \texttt{ConcreteModel\_impl}
maintient son propre état d'appartenance indépendamment des autres
modèles, ce qui est cohérent avec la sémantique de
\texttt{isInDefinition}~: un même \texttt{ConcreteObject} peut appartenir
à plusieurs \texttt{ConcreteModel} simultanément
(comportement par défaut lorsque la propriété
\texttt{canNotBeShared} de \texttt{MetaDef} n'est pas activée~;
cf.~section~\ref{sec:metal2}).

%\paragraph{Impact sur la complexité.}
Avant l'optimisation, chaque appel à \texttt{\_isInDefinition} coûtait
$O(n)$ où $n$ est le nombre de \texttt{K\_Property} dans le range du
\texttt{ConcreteModel}. Après l'optimisation, le coût est $O(1)$ pour
toute requête déjà vue, et $O(n)$ uniquement lors du premier scan pour
chaque paire (\texttt{ConcreteModel}, \texttt{ConcreteObject}). Dans le
contexte de \texttt{\_getName}, qui interroge les mêmes propriétés de
nom de manière répétitive lors du bootstrapping, le nombre de scans
effectifs tend vers $O(k)$ où $k$ est le nombre de paires distinctes
plutôt que vers $O(n \times \text{appels})$.

% -------------------------------------------------------------------
\subsection{Tests de performance --- itération~2}
\label{sec:scenario1-tests2}
% -------------------------------------------------------------------

\subsubsection{Bilan du second cycle d'optimisation}

Le tableau~\ref{tab:bilan-opti2} synthétise les gains cumulés à l'issue
des deux cycles d'optimisation.

\begin{table}[h]
\centering
\caption{Bilan cumulé des deux cycles d'optimisation ---
scénario d'initialisation. Les médianes sont calculées sur 10~répétitions
(mesures complètes en annexe~\ref{ann:mesures-s1-a}).}
\label{tab:bilan-opti2}
\begin{tabular}{lrrr}
\hline
\textbf{Indicateur} & \textbf{Baseline} & \textbf{Après~1--3} &
\textbf{Après~1--4} \\
\hline
Temps de démarrage (médiane, ms) & 18\,198 & 4\,516 & 1\,490 \\
Facteur cumulé & --- & $\times$\,4,0 & $\times$\,12,2 \\
Complexité (appartenance, niveau méta) & $O(n^2 \cdot s)$ & $O(n^2)$ & $O(n^2)$ \\
Complexité (appartenance, niveau concret) & $O(n^2)$ & $O(n^2)$ &
$O(n \cdot k)$ \\
\hline
\end{tabular}
\end{table}

Le passage de $\times$\,4,0 à $\times$\,12,2 illustre l'effet
multiplicateur d'une optimisation portant sur le coût unitaire dans une
boucle déjà serrée par les cycles précédents~: une fois la redondance
structurelle éliminée (iter.~1), même un cache à coût quasi-nul produit
un gain important.

% -------------------------------------------------------------------
\subsection{Synthèse du scénario~1}
\label{sec:scenario1-synthese}
% -------------------------------------------------------------------

Le scénario d'initialisation a permis d'identifier et de résoudre deux
familles de goulots d'étranglement successives, toutes deux enracinées
dans la même propriété architecturale de MetaDONE~: l'encodage uniforme
de l'appartenance et du typage dans un graphe homogène de
\texttt{DataObject}.

La première famille (optimisations~1--3) portait sur le niveau
méta-modèle (M2/M3) de MetaL$_2$~: le cache de \texttt{MetaModel\_impl}
était structurellement inopérant, transformant une vérification
d'appartenance $O(1)$ en un scan $O(n)$ répété, lui-même amplifié d'un
facteur~57 par la boucle de \texttt{getMOByName}. La complexité
effective du bootstrapping était $O(n^3)$.

La seconde famille (optimisation~4) portait sur le niveau concret
(M1) de MetaL$_2$~: la symétrie réflexive entre les deux niveaux
s'est retrouvée dans le problème lui-même.
\texttt{\_isInDefinition} dans \texttt{ConcreteModel\_impl}
reproduisait exactement le même antipattern, avec un facteur
d'amplification de~43.

Ces deux familles partagent une cause commune~: l'absence de mémoïsation
sur des résultats stables. Les \texttt{DataObject} d'un projet MetaDONE
sont majoritairement en lecture une fois le bootstrapping achevé~; le
coût de maintien des caches (invalidation sur mutation, mise à jour
incrémentale sur ajout) est donc faible comparé au gain sur les accès.
Cette observation est cohérente avec les recommandations générales de la
littérature sur les applications à forte intensité de lecture sur des
graphes~\cite{oaks2014java, fowler2002patterns}.

Les quatre optimisations sont ciblées, non invasives et n'affectent
pas les contrats publics des interfaces de MetaL$_2$ (aux deux
niveaux méta et concret)~: elles opèrent exclusivement dans les
classes \texttt{*\_impl} du bundle \texttt{metadone\_bundle\_metabusiness},
sans modifier le bundle \texttt{repository} qui implémente MetaL$_1$,
ni les interfaces exposées aux plugins.

\paragraph{Confrontation à l'hypothèse initiale.}
L'hypothèse formulée en méthodologie (section~\ref{sec:scenarios})
selon laquelle cette phase d'initialisation souffrirait du problème
du N+1 et des limites du \emph{lazy loading} n'a pas été validée par
les mesures~: les goulots identifiés (familles~1 et~2 ci-dessus) sont
strictement CPU-bound et algorithmiques. Aucune des chaînes d'appels
critiques ne déclenche d'accès supplémentaires à la base de données
au-delà du chargement initial~; les boucles linéaires problématiques
opèrent toutes sur des collections JPA déjà chargées dans l'identity
map. Le surcoût JPA effectif, mesurable mais minoritaire après
optimisations (\texttt{pcGetdeleted} et \texttt{getRangeOf} ci-dessous),
n'affecte pas la conclusion principale.

\paragraph{Limites et goulots résiduels.}
\texttt{pcGetdeleted} et \texttt{getRangeOf} conservent une présence
non négligeable dans les profils post-optimisation. Ces méthodes
appartiennent au bundle \texttt{metadone\_bundle\_repository\_jpa2} et
reflètent le coût intrinsèque des accès à la base de données via OpenJPA,
indépendamment des caches applicatifs. Leur réduction nécessiterait
d'agir à un niveau différent~: stratégie de chargement JPA (batch
fetching, eager vs. lazy), indexation de la colonne \texttt{range} dans
le schéma relationnel, ou remplacement partiel des requêtes par des
structures en mémoire. Ces pistes constituent les axes d'investigation
des scénarios suivants.